\section*{第 \ref{ch:g6:how-plants-spread} 章——植物如何传播}

\begin{solution}{exo:g6:how-plants-spread:1}
\omterm{prop:g6:how-plants-spread:travel}{定殖}是指一个\omterm{def:g5:biodiversity:species}{物种}占据新的地域。
\omterm{def:g6:how-plants-spread:dispersal}{扩散}是指将\omterm{def:g2:from-seed-to-plant:seed}{种子}、\omterm{prop:g3:flowers-fruits-seeds:transform}{果实}或\omterm{def:g4:plants-without-seeds:spore}{孢子}带离
母株；携带者有风、\omterm{def:g1:animals-around-us:animal}{动物}（体外和体内）、
水——以及一些能自行弹射的\omterm{prop:g3:flowers-fruits-seeds:transform}{果实}。
\end{solution}

\begin{solution}{exo:g6:how-plants-spread:2}
若\omterm{def:g2:from-seed-to-plant:seed}{种子}落在母株脚下，它就会在母株的荫蔽下\omterm{def:g2:from-seed-to-plant:germination}{萌发}，
并与母株的根争夺光照和水分，
而对手的体积是它的一百倍。
距离使它避开了这场注定失败的竞争。
\end{solution}

\begin{solution}{exo:g6:how-plants-spread:3}
蒲公英的冠毛：靠风。牛蒡的刺球：靠体外的\omterm{def:g1:animals-around-us:animal}{动物}毛皮。
樱桃：靠体内的\omterm{def:g1:animals-around-us:animal}{动物}\omterm{def:g4:journey-of-food:tract}{消化道}。枫树的翅果：靠风（旋转的翅）。
椰子：靠水。
\end{solution}

\begin{solution}{exo:g6:how-plants-spread:4}
它们坚韧的种皮使其能够安然通过\omterm{prop:g3:sorting-living-things:groups}{鸟类}的\omterm{def:g4:journey-of-food:tract}{消化道}，
\omterm{def:g2:from-seed-to-plant:seed}{种子}随鸟粪落在远离母株的地方，
并附带一定剂量的肥料。
\end{solution}

\begin{solution}{exo:g6:how-plants-spread:5}
自行弹射：荚果干燥后张紧，如同弹簧，突然裂开，
将\omterm{def:g2:from-seed-to-plant:seed}{种子}射到离母株数米远的地方。
\end{solution}

\begin{solution}{exo:g6:how-plants-spread:6}
偃麦草借地下茎向前蔓延；悬钩子的枝条弯向地面，
在顶端生根（或草莓借\omterm{ex:g4:plants-without-seeds:runner}{匍匐茎}
跳跃式扩展）。
\end{solution}

\begin{solution}{exo:g6:how-plants-spread:7}
蒲公英：借冠毛乘风飞来。罂粟：\omterm{def:g2:from-seed-to-plant:seed}{种子}早已埋在土壤中，
土地被翻动后便被唤醒。
黑莓：由停在枝头的\omterm{prop:g3:sorting-living-things:groups}{鸟类}带来，\omterm{def:g2:from-seed-to-plant:seed}{种子}藏在
粪便里。牛蒡：挂在狗的毛皮上被带来。
\end{solution}

\begin{solution}{exo:g6:how-plants-spread:8}
(a) 飞行装置——冠毛：靠风。(b) 果肉鲜艳、
种皮坚韧的\omterm{def:g2:from-seed-to-plant:seed}{种子}在体内：被\omterm{def:g1:animals-around-us:animal}{动物}吃掉，\omterm{def:g2:from-seed-to-plant:seed}{种子}随
\omterm{def:g4:journey-of-food:tract}{消化道}转移。(c) 能勾住羊毛的钩刺：靠\omterm{def:g1:animals-around-us:animal}{动物}毛皮——
而那件毛衣的实验恰好证实了这一点。
\end{solution}

\begin{solution}{exo:g6:how-plants-spread:9}
顶端生根、弯向地面的茎完成的是“米”级的扩展——
在邻近的地面上密集而稳妥地占据。被鸟吃掉的浆果
完成的是“英里”级的扩展——偶尔的远距离落地，
在远处找到了新的落脚点。
\end{solution}

\begin{solution}{exo:g6:how-plants-spread:10}
\omterm{def:g4:plants-without-seeds:spore}{孢子}能到达两岸——它们到处漂浮——但\omterm{prop:g4:plants-without-seeds:damp}{蕨类}
的繁殖需要潮湿，因此只有北岸的\omterm{def:g6:exploring-our-environment:conditions}{条件}接纳了它们。
\omterm{def:g6:how-plants-spread:dispersal}{扩散}负责送达，\omterm{def:g6:exploring-our-environment:conditions}{条件}决定谁能留下：
居住者与\omterm{def:g6:exploring-our-environment:conditions}{条件}相匹配，
无论有多少外来者到达。
\end{solution}

\begin{solution}{exo:g6:how-plants-spread:11}
只有远距离的传播方式：风（有冠毛、极细的\omterm{def:g2:from-seed-to-plant:seed}{种子}和\omterm{def:g4:plants-without-seeds:spore}{孢子}），
海水（能漂浮的\omterm{prop:g3:flowers-fruits-seeds:transform}{果实}，如椰子），
以及\omterm{prop:g3:sorting-living-things:groups}{鸟类}（肉质\omterm{prop:g3:flowers-fruits-seeds:transform}{果实}的\omterm{def:g2:from-seed-to-plant:seed}{种子}经\omterm{def:g4:journey-of-food:tract}{消化道}传播）。由此可以预测，
先锋\omterm{def:g1:plants-around-us:plant}{植物}应具有冠毛、极细、能漂浮或借浆果传播的装置——
而不是沉重且无运载者的\omterm{def:g2:from-seed-to-plant:seed}{种子}。
\end{solution}

\begin{solution}{exo:g6:how-plants-spread:12}
糖分是为运输付出的报酬：\omterm{prop:g3:sorting-living-things:groups}{鸟类}以果肉作代价，
把\omterm{def:g2:from-seed-to-plant:seed}{种子}带离母株的荫蔽，
并连同肥料放下——这是距离，即第一份礼物。坚韧的种皮则是另一半，
因为运送者是\omterm{def:g4:journey-of-food:tract}{消化道}：报酬必须能被\omterm{def:g4:journey-of-food:digestion}{消化}，
而乘客则不能被\omterm{def:g4:journey-of-food:digestion}{消化}。
甜味招募运送者，种皮则使\omterm{def:g2:from-seed-to-plant:seed}{种子}在其体内存活。
\end{solution}

\begin{solution}{pb:g6:how-plants-spread:1}
\textbf{1.} 罂粟：\omterm{def:g2:from-seed-to-plant:seed}{种子}早已在被耕翻的土壤中，
土地翻动后便被唤醒——这是\omterm{prop:g6:life-through-seasons:herbs}{一年生植物}储备的\omterm{def:g2:from-seed-to-plant:seed}{种子}。蒲公英和
蓟：借有冠毛的\omterm{def:g2:from-seed-to-plant:seed}{种子}乘风飞来。

\textbf{2.} 在别处吃了黑莓的\omterm{prop:g3:sorting-living-things:groups}{鸟类}停在栅栏柱上，
把\omterm{def:g2:from-seed-to-plant:seed}{种子}连同肥料排在柱子下方——
因此\omterm{prop:g2:from-seed-to-plant:cycle}{幼苗}标示的是栖木的位置，而不是风的路径。

\textbf{3.} 原因在于干燥：开阔向阳的田野缺少
\omterm{prop:g4:plants-without-seeds:damp}{蕨类}娇嫩的\omterm{def:g4:plants-without-seeds:spore}{孢子}发育阶段所需的潮湿。
\omterm{def:g4:plants-without-seeds:spore}{孢子}已经到达，但未能\omterm{prop:g6:how-plants-spread:travel}{定殖}。

\textbf{4.} \omterm{def:g4:plants-without-seeds:vegetative}{营养繁殖}式的扩展——借地下
茎蔓延。它只能从已\omterm{prop:g6:how-plants-spread:travel}{定殖}的植株向邻近地面扩展，
因此必须从边缘逐步进入；
它没有可在中央着陆的空中阶段。

\textbf{5.} 以\omterm{def:g2:from-seed-to-plant:seed}{种子}的形式留在土壤中：在长期竞争中
落败的\omterm{prop:g6:life-through-seasons:herbs}{一年生植物}以\omterm{def:g2:from-seed-to-plant:seed}{种子}库的形式留存下来，
等待下一块裸地。

\textbf{6.} 体内的\omterm{def:g1:animals-around-us:animal}{动物}传播（浆果被吃掉，
\omterm{def:g2:from-seed-to-plant:seed}{种子}从栖木上排出——在篱笆线和树篱下最为密集），
对山楂同样如此，同样由\omterm{prop:g3:sorting-living-things:groups}{鸟类}运送；
栖木式的分布格局便是证据。

\textbf{7.} 橡子和栗子很重，无人帮助便无法远行——
它们落在母树脚下。其运送者是松鸦
（以及松鼠）：它们把橡子埋作冬季储备，
并遗忘其中一部分，于是在
远离母树的地方长出橡树。

\textbf{8.} 每年长出的植被遮蔽其下的地面，
使地表变凉、变暗：喜阳的先锋\omterm{def:g1:plants-around-us:plant}{植物}
自身的成功却消除了它们赖以生存的\omterm{def:g6:exploring-our-environment:conditions}{条件}，
而改变后的\omterm{def:g6:exploring-our-environment:conditions}{条件}接着有利于耐阴\omterm{def:g1:plants-around-us:plant}{植物}——
定居者改变了场地，被改变的场地又选择新的定居者。

\textbf{9.} 第一波，土壤中等待的\omterm{def:g2:from-seed-to-plant:seed}{种子}（瞬时——它们早已在那里）。
第二波，借风传播的\omterm{def:g2:from-seed-to-plant:seed}{种子}（数周：
空气持续输送）。第三波，由\omterm{def:g1:animals-around-us:animal}{动物}携带的灌木
（数年：逐根栖木接力）。第四波，\omterm{def:g2:from-seed-to-plant:seed}{种子}大而重的乔木（最晚：
它们要等待松鸦和松鼠埋下又遗忘的\omterm{def:g2:from-seed-to-plant:seed}{种子}）。

\textbf{10.} 缝隙里没有容身之地：灌木的根需要
路面接缝所缺少的土壤体积；而演替的
后续各波还需要它们的运送者和\omterm{def:g6:exploring-our-environment:conditions}{条件}——栖息\omterm{prop:g3:sorting-living-things:groups}{鸟类}的
粪便大多落不进缝隙，被烤干的接缝拒绝了
所有真正落下的\omterm{def:g2:from-seed-to-plant:seed}{种子}。第一年借风传播的\omterm{def:g2:from-seed-to-plant:seed}{种子}，
是这片场地唯一能接受的运送。

\textbf{11.} 演替的趋势是：裸地、草本、灌木、
先锋乔木——每一波都更高大、更荫蔽：延续下去，
最终形成林地，柳树和桦树已经表明了这一点。
例外是割草或放牧的草地，那里人为干预每年
都会清除灌木一波——这证明只有
\emph{任其自然}，场地才会走向森林。

\textbf{12.} 例如：“\omterm{def:g6:how-plants-spread:dispersal}{扩散}将定居者送到裸地——来自土壤、
风和\omterm{prop:g3:sorting-living-things:groups}{鸟类}——
并且在每个阶段，场地的\omterm{def:g6:exploring-our-environment:conditions}{条件}决定了哪些到来者能够留下。
定居因而分波次进行，每一波都改变\omterm{def:g6:exploring-our-environment:conditions}{条件}，
并选择下一波。”
\end{solution}
